% Template Laporan Proyek 2 D4 Teknik Informatika ULBI
% Sesuai Petunjuk Teknis ll.my.id/instruction (Capstone Project Vokasi)

\documentclass[12pt,a4paper]{article}

% Packages
\usepackage[utf8]{inputenc}
\usepackage[T1]{fontenc}
\usepackage{mathptmx} % Times New Roman
\usepackage{courier}  % Courier for monospace
\usepackage{helvet}   % Helvetica for sans-serif
\usepackage{amsmath}
\usepackage[english,indonesian]{babel}
\usepackage{geometry}
\usepackage{setspace}
\usepackage{titlesec}
\usepackage{graphicx}
\usepackage{hyperref}
\usepackage{booktabs}
\usepackage{array}
\usepackage[format=hang,font=footnotesize,labelfont=bf]{caption}
\usepackage{enumitem}
\usepackage{parskip}
\usepackage{fancyhdr}
\usepackage{tocloft}
\usepackage{tabularx}
\usepackage{longtable}
\usepackage{ragged2e}

% Page Setup
\geometry{
    top=3cm,
    bottom=3cm,
    left=4cm,
    right=3cm
}

% Title Formatting
\titleformat{\section}
    {\normalfont\fontsize{14}{17}\bfseries\centering}{\thesection}{1em}{}
    
\titleformat{\subsection}
    {\normalfont\fontsize{12}{14}\bfseries}{\thesubsection}{1em}{}

\titleformat{\subsubsection}
    {\normalfont\fontsize{12}{14}\normalfont\bfseries}{\thesubsubsection}{1em}{}

% Spacing
\onehalfspacing

% Header/Footer
\pagestyle{fancy}
\fancyhf{}
\fancyfoot[C]{\thepage}
\renewcommand{\headrulewidth}{0pt}

% Hyperref setup
\hypersetup{
    colorlinks=true,
    linkcolor=black,
    citecolor=black,
    filecolor=black,
    urlcolor=black,
}

% Table of Contents customization
\renewcommand{\contentsname}{DAFTAR ISI}
\renewcommand{\cfttoctitlefont}{\hfill\fontsize{14}{17}\selectfont\bfseries}
\renewcommand{\cftaftertoctitle}{\hfill}

% List of Tables customization
\renewcommand{\listtablename}{DAFTAR TABEL}
\renewcommand{\cftlottitlefont}{\hfill\fontsize{14}{17}\selectfont\bfseries}
\renewcommand{\cftafterlottitle}{\hfill}

% List of Figures customization
\renewcommand{\listfigurename}{DAFTAR GAMBAR}
\renewcommand{\cftloftitlefont}{\hfill\fontsize{14}{17}\selectfont\bfseries}
\renewcommand{\cftafterloftitle}{\hfill}

% Subsection numbering format
\renewcommand{\thesubsection}{\arabic{section}.\arabic{subsection}}

% Bibliography Style - APA-like with BibTeX
\usepackage[
    backend=bibtex,
    style=authoryear,
    sorting=nyt,
    maxbibnames=99,
    dashed=false,
    url=false,
    doi=true,
]{biblatex}

\urlstyle{same}
\addbibresource{references.bib}

% Make \cite behave like \parencite
\let\cite\parencite

\renewcommand*{\finalnamedelim}{\addspace\&\addspace}
\renewcommand*{\nameyeardelim}{\addcomma\space}
\renewbibmacro{in:}{}
\DeclareFieldFormat{pages}{\normalfont #1}
\DeclareFieldFormat{booktitle}{\addspace\textit{#1}}
\DeclareFieldFormat{journaltitle}{\addspace\textit{#1}}

\setcounter{biburllcpenalty}{7000}
\setcounter{biburlucpenalty}{8000}
\emergencystretch=1em
\appto{\bibsetup}{\sloppy\RaggedRight}

% Document Metadata
\newcommand{\docTitle}{RANCANG BANGUN CROOTAGENT: AGEN KECERDASAN BUATAN BERBASIS CLI ULTRA-RINGAN DENGAN MULTI-API SESSION DAN ORKESTRASI PERKAKAS PENGEMBANG BERBASIS RUST}
\newcommand{\docName}{Ahmad Zakky Nurfauzan}
\newcommand{\docNPM}{714250028}
\newcommand{\docYear}{2026}
\newcommand{\docPembimbing}{Rolly Maulana Awangga, S.T., M.T., CAIP.}

\begin{document}

% Force TOC title to uppercase
\renewcommand{\contentsname}{DAFTAR ISI}
\renewcommand{\listtablename}{DAFTAR TABEL}
\renewcommand{\listfigurename}{DAFTAR GAMBAR}

\rmfamily

% ==================== HALAMAN JUDUL ====================
\begin{titlepage}
    \centering
    \vspace*{0.5cm}
    
    {\fontsize{14}{18}\selectfont
    \textbf{LAPORAN PROYEK 2}\par}
    \vspace{0.8cm}
    
    {\fontsize{13}{16}\selectfont
    \textbf{\docTitle}\par}
    
    \vspace{1.2cm}
    
    {\fontsize{12}{14}\selectfont
    Diajukan sebagai salah satu syarat untuk kelulusan mata kuliah Proyek 2\par
    pada Program Studi Sarjana Terapan Teknik Informatika\par}
    
    \vspace{1.5cm}
    
    {\fontsize{12}{14}\selectfont
    Oleh:\par
    \textbf{\docName}\par
    \textbf{NPM: \docNPM}\par}
    
    \vspace{1.5cm}
    
    \IfFileExists{logo.png}{
        \includegraphics[width=0.35\textwidth]{logo.png}
    }{
        {\fontsize{14}{16}\bfseries [LOGO ULBI]\par}
    }
    
    \vfill
    
    {\fontsize{12}{15}\selectfont
    \textbf{PROGRAM STUDI SARJANA TERAPAN TEKNIK INFORMATIKA}\par
    \textbf{SEKOLAH TINGGI / FAKULTAS LOGISTIK DAN BISNIS}\par
    \textbf{UNIVERSITAS LOGISTIK DAN BISNIS INTERNASIONAL}\par
    \textbf{BANDUNG}\par
    \textbf{\docYear}\par}
    
\end{titlepage}
\newpage

% ==================== HALAMAN PENGESAHAN ====================
\begin{center}
    {\fontsize{14}{17}\selectfont\textbf{HALAMAN PENGESAHAN}\par}
    \vspace{1.2cm}
    
    {\fontsize{12}{15}\selectfont
    \textbf{\docTitle}\par}
    
    \vspace{1cm}
    
    Oleh:\par
    \textbf{\docName}\par
    \textbf{NPM: \docNPM}\par
    
    \vspace{1cm}
    Laporan Proyek 2 ini telah disetujui dan disahkan pada:\par
    Bandung, 17 November 2026\par
    
    \vspace{1.5cm}
    
    Menyetujui,\par
    Dosen Pembimbing,\par
    
    \vspace{2.5cm}
    
    \textbf{\underline{\docPembimbing}}\par
    NIDN: 0423127801\par
    
    \vspace{1.5cm}
    
    Mengetahui,\par
    Ketua Program Studi Sarjana Terapan Teknik Informatika,\par
    
    \vspace{2.5cm}
    
    \textbf{\underline{M. Yusril Helmi Setyawan, S.Kom., M.Kom.}}\par
    NIDN: 0425098901\par
\end{center}
\newpage

% ==================== DAFTAR ISI DLL (Angka Romawi) ====================
\pagenumbering{roman}
\setcounter{page}{1}

% RINGKASAN
\section*{RINGKASAN}
\phantomsection
\addcontentsline{toc}{section}{RINGKASAN}
Perkembangan asisten pengembang berbasis \textit{Large Language Model} (LLM) seperti Claude Code, OpenAI Codex CLI, Google Antigravity CLI, dan AWS Kiro telah merevolusi produktivitas rekayasa perangkat lunak. Namun, sebagian besar kakas agen tersebut diimplementasikan menggunakan lingkungan terinterpretasi seperti Node.js atau Python yang memakan sumber daya memori besar (rata-rata 150--300 MB RAM saat \textit{idle}) serta membutuhkan \textit{runtime} eksternal yang rumit untuk didistribusikan ke lintas sistem operasi.

Pada Proyek 2 ini, dibangun sistem agen cerdas berbasis CLI bernama \textbf{CRootAgent} yang dikembangkan menggunakan bahasa pemrograman sistem \textbf{Rust}. CRootAgent memiliki arsitektur \textit{native} ultra-ringan dengan konsumsi memori dasar hanya 10,2 MB RAM dan waktu \textit{startup} di bawah 20 milidetik. Sistem ini mengintegrasikan protokol \textit{multi-API session} yang mendukung lima penyedia kecerdasan buatan terkemuka (Google Gemini, Anthropic Claude, OpenAI, Groq Cloud, dan Ollama lokal) dengan kemampuan alih-sesi dinamis. Selain itu, CRootAgent mengimplementasikan siklus penalaran ReAct (\textit{Reasoning and Acting}) yang dilengkapi lima perkakas operasi berkas dan terminal lokal (\texttt{read\_file}, \texttt{write\_file}, \texttt{edit\_file}, \texttt{run\_command}, dan \texttt{list\_dir}) dengan menerapkan mekanisme keamanan \textit{Human-in-the-Loop} (HITL).

Seluruh luaran yang ditargetkan pada Proyek 2 telah berhasil diselesaikan dan diverifikasi:
\begin{enumerate}[noitemsep]
    \item Pembangunan inti mesin agen CLI berbasis Rust (\texttt{CRootAgent/crootagent}) dengan 10 kasus uji otomatis (\textit{unit} dan \textit{integration test}) dengan tingkat kelulusan 100\%.
    \item Pipeline CI/CD GitHub Actions yang secara otomatis menjalankan pengujian, melakukan \textit{auto-tagging}, dan membangun \textit{binary release} mandiri untuk empat target sistem operasi: Linux x86\_64, macOS Apple Silicon (arm64), macOS Intel (x86\_64), dan Windows x86\_64.
    \item Pembuatan portal dokumentasi dan landing page web (\texttt{crootagent.github.io}) yang mematuhi konvensi CrootJS serta memanggil API GitHub Releases secara langsung untuk mendeteksi sistem operasi pengguna dan menyajikan tombol unduh 1-klik beserta skrip instalasi instan (\texttt{install.sh} dan \texttt{install.ps1}).
    \item Penyusunan materi presentasi spasial 3D berbasis CSS murni (\texttt{crootss.if.co.id}) dan laporan akademik LaTeX ini sesuai standar integritas petunjuk proyek vokasi.
\end{enumerate}

\textbf{Kata Kunci:} Agen CLI, Rust, Multi-API Session, ReAct Pattern, Human-in-the-Loop, CrootJS, CI/CD Multi-OS.

\clearpage

% DAFTAR ISI
\phantomsection
\addcontentsline{toc}{section}{DAFTAR ISI}
\tableofcontents
\clearpage

% DAFTAR TABEL
\phantomsection
\addcontentsline{toc}{section}{DAFTAR TABEL}
\listoftables
\clearpage

% DAFTAR GAMBAR
\phantomsection
\addcontentsline{toc}{section}{DAFTAR GAMBAR}
\listoffigures
\clearpage

% ==================== ISI LAPORAN (Angka Arab) ====================
\setcounter{section}{0}
\pagenumbering{arabic}
\setcounter{page}{1}

% BAB 1 - PENDAHULUAN
\section*{BAB 1 \\ PENDAHULUAN}
\phantomsection
\addcontentsline{toc}{section}{BAB 1 PENDAHULUAN}
\refstepcounter{section}
\subsection{Latar Belakang}
Pemanfaatan kecerdasan buatan berbasis model bahasa besar (\textit{Large Language Model} / LLM) telah mengalami pergeseran paradigma dari antarmuka obrolan berbasis web menuju agen cerdas (\textit{agentic AI}) yang beroperasi langsung di lingkungan kerja lokal pengembang perangkat lunak. Kakas seperti Claude Code dari Anthropic, OpenAI Codex CLI, Google Antigravity CLI, dan AWS Kiro membuktikan bahwa agen yang mampu membaca berkas, menulis kode, memverifikasi galat sintaks, dan mengeksekusi perintah terminal mampu meningkatkan produktivitas rekayasa perangkat lunak secara signifikan \cite{hou2024llm4se}.

Meskipun demikian, sebagian besar kakas agen yang beredar saat ini dibangun di atas lingkungan terinterpretasi seperti Node.js atau Python. Ketergantungan tersebut menimbulkan dua kendala teknis utama:
\begin{enumerate}
    \item \textbf{Konsumsi Memori dan Beban Runtime Tinggi:} Lingkungan Node.js dan Python memerlukan \textit{runtime engine} (V8 atau Python VM) yang menyerap memori dasar sebesar 150--300 MB RAM hanya untuk proses \textit{idle}, dengan waktu \textit{startup} berkisar 200--600 milidetik. Pada perangkat dengan spesifikasi terbatas atau ketika pengembang menjalankan puluhan kontainer dan proses kompilasi secara simultan, beban memori ini menjadi penghambat kinerja yang nyata \cite{pereira2021energy}.
    \item \textbf{Kerumitan Distribusi Multi-OS:} Distribusi aplikasi agen berbasis Node.js/Python menuntut pengguna menginstal paket pengelola (\texttt{npm}, \texttt{pip}), menangani konflik versi pustaka, atau mengunduh bundel biner besar (> 80 MB) yang membungkus seluruh runtime.
\end{enumerate}

Bahasa pemrograman sistem **Rust** menawarkan solusi ideal terhadap kedua kendala di atas. Rust mengompilasi kode sumber langsung menjadi instruksi mesin asli (\textit{native machine code}) tanpa \textit{garbage collector} atau \textit{virtual machine}, menjamin keamanan memori waktu-kompilasi (\textit{compile-time memory safety}), dan menghasilkan biner mandiri yang sangat ringkas dengan konsumsi memori hanya berkisar 8--15 MB RAM \cite{levy2020rust}.

Selain efisiensi performa, integrasi multi-penyedia kecerdasan buatan (\textit{multi-API provider}) menjadi kebutuhan penting. Pengembang kerap perlu beralih antara Google Gemini (yang menawarkan jendela konteks gratis hingga 1 juta token), Groq Cloud (dengan latensi pemrosesan Llama 3.3 yang sangat tinggi), OpenAI, Anthropic Claude, hingga Ollama lokal untuk pekerjaan tanpa koneksi internet.

Berdasarkan latar belakang tersebut, pada Proyek 2 Program Studi Sarjana Terapan Teknik Informatika Universitas Logistik dan Bisnis Internasional (ULBI), dirancang dan dibangun sistem agen cerdas berbasis CLI bernama \textbf{CRootAgent}. Proyek ini sekaligus menjadi pembuktian capaian profil kompetensi **Software Engineer** (Proyek 2), di mana mahasiswa dituntut membangun fitur lengkap \textit{end-to-end} mulai dari analisis kebutuhan, perancangan arsitektur, implementasi kode sumber yang modular, hingga pengujian otomatis yang terintegrasi pada pipeline CI/CD.

\subsection{Rumusan Masalah}
Berdasarkan latar belakang di atas, rumusan masalah dalam Proyek 2 ini adalah:
\begin{enumerate}
    \item Bagaimana merancang dan mengimplementasikan agen pengembang CLI yang ultra-ringan dan berkinerja tinggi menggunakan bahasa pemrograman Rust?
    \item Bagaimana membangun arsitektur \textit{multi-API session} yang mampu menghubungkan berbagai penyedia LLM (OpenAI, Anthropic Claude, Google Gemini, Groq, dan Ollama) dengan manajemen sesi yang persisten?
    \item Bagaimana mengimplementasikan siklus ReAct yang dilengkapi mekanisme keselamatan \textit{Human-in-the-Loop} (HITL) saat mengeksekusi perkakas lokal?
    \item Bagaimana merancang pipeline otomatisasi CI/CD di GitHub Actions untuk pengujian kode dan rilis kompilasi multi-OS (Linux, macOS, dan Windows) secara mandiri?
    \item Bagaimana membangun portal web \textit{landing page} yang mematuhi konvensi CrootJS untuk mendeteksi rilis biner terbaru secara otomatis?
\end{enumerate}

\subsection{Tujuan Proyek}
Tujuan yang hendak dicapai melalui pelaksanaan Proyek 2 ini adalah:
\begin{enumerate}
    \item Mengembangkan aplikasi biner mandiri \textbf{CRootAgent} menggunakan Rust dengan konsumsi memori di bawah 20 MB dan waktu \textit{startup} di bawah 50 milidetik.
    \item Mengimplementasikan modul perutean \textit{multi-API} yang modular dan mendukung alih penyedia/model secara dinamis.
    \item Menyediakan modul perkakas pengembang lokal (\texttt{read\_file}, \texttt{write\_file}, \texttt{edit\_file}, \texttt{run\_command}, dan \texttt{list\_dir}) yang aman dan transparan.
    \item Membangun alur CI/CD otomatis di GitHub Actions yang melakukan pengujian unit/integrasi, pemberian tag versi otomatis (\textit{auto-tagging}), serta pembuatan rilis biner multi-OS.
    \item Membangun portal web \texttt{crootagent.github.io} menggunakan CrootJS tanpa dependensi berat untuk memfasilitasi distribusi perangkat lunak.
\end{enumerate}

\subsection{Batasan Masalah}
Batasan masalah dalam pelaksanaan Proyek 2 ini meliputi:
\begin{enumerate}
    \item Antarmuka pengguna berfokus pada baris perintah (\textit{Command Line Interface} / CLI) dan terminal interaktif (REPL), tidak mencakup antarmuka grafis desktop (GUI).
    \item Penyedia LLM yang didukung secara resmi meliputi Google Gemini, Anthropic Claude, OpenAI, Groq Cloud, dan Ollama lokal.
    \item Pengujian otomatis difokuskan pada uji unit konfigurasi, manajemen sesi, dan integritas eksekusi perkakas lokal.
    \item Target sistem operasi pada pipeline rilis otomatis dibatasi pada Linux x86\_64, macOS Apple Silicon (arm64), macOS Intel (x86\_64), dan Windows x86\_64.
\end{enumerate}

\subsection{Manfaat Proyek}
Manfaat dari pengembangan proyek ini mencakup:
\begin{enumerate}
    \item \textbf{Bagi Pengembang Perangkat Lunak:} Menyediakan kakas asisten pengembang yang sangat cepat, hemat daya baterai dan memori, serta bebas digunakan dengan berbagai kunci API pilihan sendiri.
    \item \textbf{Bagi Institusi (ULBI):} Menghasilkan produk perangkat lunak bernilai guna tinggi yang menjadi bukti nyata kompetensi rekayasa perangkat lunak vokasi tingkat Sarjana Terapan.
    \item \textbf{Bagi Mahasiswa:} Mengasah penguasaan bahasa pemrograman sistem Rust, perancangan arsitektur \textit{agentic AI}, manajemen SCM GitHub, dan otomatisasi CI/CD modern.
\end{enumerate}

\subsection{Profil Kompetensi Capstone Proyek 2}
Sesuai dengan ketentuan petunjuk proyek capstone bertahap di lingkungan pendidikan vokasi \cite{awangga2024agentic}, Proyek 2 menuntut pembuktian profil **Software Engineer**:
\begin{itemize}
    \item Membangun fitur lengkap \textit{end-to-end}: analisis kebutuhan $\rightarrow$ desain sistem $\rightarrow$ implementasi $\rightarrow$ integrasi.
    \item Pengujian yang ditulis sendiri oleh mahasiswa: uji unit dan integrasi dengan verifikasi mekanis nyata.
    \item Pipeline CI/CD yang secara otomatis menjalankan \textit{test suite} setiap kali terjadi perubahan kode sumber.
\end{itemize}
Seluruh poin kompetensi tersebut dibuktikan dan dilaporkan secara transparan pada dokumen ini.

\clearpage

% BAB 2 - TINJAUAN PUSTAKA
\section*{BAB 2 \\ TINJAUAN PUSTAKA DAN KAJIAN TERKAIT}
\phantomsection
\addcontentsline{toc}{section}{BAB 2 TINJAUAN PUSTAKA DAN KAJIAN TERKAIT}
\refstepcounter{section}
\subsection{Arsitektur Agen Cerdas dan Pola ReAct}
Sistem agen cerdas berbasis model bahasa besar tidak sekadar berperan sebagai model generasi teks pasif, melainkan bertindak sebagai pengambil keputusan yang mampu berinteraksi dengan lingkungannya. Pola arsitektur \textbf{ReAct} (\textit{Reasoning and Acting}) yang diperkenalkan oleh \cite{yao2023react} menggabungkan jejak penalaran verbal (\textit{reasoning traces}) dengan tindakan konkret (\textit{actions}). 

Dalam siklus ReAct, agen menerima masukan dari pengguna dan menjalankan tahapan berulang:
\begin{enumerate}
    \item \textbf{Reason (Pemikiran):} Model mengevaluasi kondisi konteks percakapan saat ini dan merumuskan sub-tujuan logis.
    \item \textbf{Act (Tindakan):} Model menerbitkan pemanggilan fungsi (\textit{tool calling}) dengan menyertakan nama perkakas dan argumen terstruktur dalam format JSON.
    \item \textbf{Observe (Pengamatan):} Lingkungan eksekusi lokal menjalankan perkakas tersebut dan mengembalikan hasilnya (\textit{observation}) ke dalam konteks pesan model.
\end{enumerate}
Pendekatan ini terbukti mengurangi halusinasi secara drastis serta memungkinkan model menyelesaikan tugas rekayasa perangkat lunak multi-langkah seperti pencarian berkas, perbaikan bug, dan kompilasi \cite{schick2023toolformer}.

\subsection{Kajian Komparatif Kakas Agen CLI Pengembang}
Tabel \ref{tab:sota} menyajikan perbandingan antara beberapa kakas agen coding CLI terkemuka di industri perangkat lunak dengan CRootAgent yang dikembangkan pada proyek ini.

\begin{table}[htbp]
\centering
\small
\caption{Perbandingan State-of-the-Art Kakas Agen Coding CLI}
\label{tab:sota}
\begin{tabularx}{\textwidth}{l p{2.8cm} p{2.2cm} p{2.2cm} X}
\toprule
\textbf{Kakas Agen} & \textbf{Bahasa / Runtime} & \textbf{Multi-API Provider} & \textbf{Jejak Memori (RAM Idle)} & \textbf{Karakteristik Utama} \\
\midrule
Claude Code & Node.js / TypeScript & Tunggal (Anthropic) & $\sim$160 MB & Kuat pada integrasi Claude, distribusi via npm. \\
OpenAI Codex CLI & Python / Node.js & Tunggal (OpenAI) & $\sim$180 MB & Berfokus pada model OpenAI, ketergantungan runtime. \\
Google Antigravity & Node.js / Go Wrapper & Multi-Provider & $\sim$140 MB & Ekosistem ekstensi luas, beban startup terpengaruh Node. \\
AWS Kiro & TypeScript & Multi-Provider & $\sim$150 MB & Steering file fleksibel, instalasi via npm. \\
\textbf{CRootAgent} & \textbf{Rust (Native)} & \textbf{Multi-Provider} & \textbf{10,2 MB} & \textbf{Binary mandiri, zero runtime dependency, multi-API session, HITL safety.} \\
\bottomrule
\end{tabularx}
\end{table}

\subsection{Karakteristik Bahasa Pemrograman Sistem Rust}
Rust dipilih sebagai fondasi pengembangan CRootAgent karena keunggulan intrinsiknya dalam rekayasa perangkat lunak sistem:
\begin{enumerate}
    \item \textbf{Keamanan Memori Tanpa Garbage Collector:} Rust menerapkan sistem kepemilikan (\textit{ownership}) dan peminjaman (\textit{borrow checker}) yang memvalidasi integritas memori pada saat kompilasi. Karakteristik ini mengeliminasi celah kerentanan seperti \textit{dangling pointer}, \textit{double free}, dan \textit{data race} tanpa membebankan penalti jeda \textit{garbage collector} \cite{levy2020rust}.
    \item \textbf{Jejak Memori Ringkas dan Efisiensi Daya:} Penelitian empiris oleh \cite{pereira2021energy} menunjukkan bahwa biner hasil kompilasi Rust menempati peringkat teratas dalam efisiensi energi dan kecepatan eksekusi, mendekati performa bahasa C/C++ namun dengan jaminan keamanan memori penuh.
    \item \textbf{Ekosistem Asinkronus Modern:} Pustaka \texttt{tokio} dan \texttt{reqwest} pada Rust menyediakan mesin runtime I/O asinkronus berbasis \textit{epoll/kqueue} yang mampu menangani ratusan koneksi jaringan dengan konsumsi memori minimal.
\end{enumerate}

\subsection{Prinsip Human-in-the-Loop (HITL) dan Explainable AI (xAI)}
Dalam sistem AI agentic, pendelegasian wewenang penuh kepada model untuk mengeksekusi perintah sistem operasi tanpa pengawasan manusia membawa risiko keselamatan fatal (\textit{destructive actions}), seperti penghapusan berkas proyek secara tidak disengaja atau eksekusi perintah terminal yang berbahaya \cite{amershi2019human}. 

Prinsip keselamatan yang diimplementasikan pada CRootAgent mengadopsi standar sistem agentic terapan:
\begin{itemize}
    \item \textbf{Human-in-the-Loop (HITL):} AI hanya bertindak sebagai pengusul tindakan (\textit{recommender}), sementara keputusan eksekusi berada di tangan manusia. Setiap perkakas yang berstatus destruktif (\texttt{write\_file}, \texttt{edit\_file}, dan \texttt{run\_command}) wajib memicu dialog konfirmasi interaktif \texttt{[y/N]} sebelum dieksekusi, kecuali bila pengguna secara sadar mengaktifkan flag bypass (\texttt{-y} / \texttt{--yes}).
    \item \textbf{Explainable AI (xAI) dan Provenance:} Setiap respon yang dihasilkan menyertakan atribusi model nyata yang memproses (\textit{provider} dan nama model asli), serta mencetak argumen dan keluaran perkakas secara transparan pada antarmuka terminal.
\end{itemize}

\subsection{Konvensi CrootJS untuk Antarmuka Web}
Untuk mendistribusikan berkas biner CRootAgent kepada pengguna, dibangun landing page web menggunakan konvensi **CrootJS** (\url{https://croot.js.org/docs/conventions}). Konvensi ini mewajibkan pemisahan berkas HTML, CSS, dan JavaScript secara independen, memuat modul langsung melalui CDN tanpa \textit{bundler} (\texttt{webpack}, \texttt{vite}), serta mengimpor modul fungsional kecil seperti \texttt{api.js} untuk panggilan asinkronus dan \texttt{useragent.js} untuk deteksi platform. Pendekatan ini selaras dengan prinsip kesederhanaan dan keandalan sistem.

\clearpage

% BAB 3 - METODOLOGI DAN PERANCANGAN SISTEM
\section*{BAB 3 \\ METODOLOGI DAN PERANCANGAN SISTEM}
\phantomsection
\addcontentsline{toc}{section}{BAB 3 METODOLOGI DAN PERANCANGAN SISTEM}
\refstepcounter{section}
\subsection{Metodologi Rekayasa Perangkat Lunak}
Pengembangan sistem CRootAgent menerapkan metodologi rekayasa perangkat lunak iteratif yang berorientasi pada kualitas dan pengujian (\textit{test-driven approach}). Tahapan metodologi meliputi:
\begin{enumerate}
    \item \textbf{Analisis Kebutuhan:} Mengidentifikasi spesifikasi fungsional dan kriteria performa agen coding CLI.
    \item \textbf{Perancangan Arsitektur:} Merancang pemisahan tanggung jawab antarmodul (\textit{Separation of Concerns}) menggunakan pola \textit{trait-based abstraction} khas Rust.
    \item \textbf{Implementasi Fitur End-to-End:} Mengembangkan pustaka inti, antarmuka baris perintah, adaptor penyedia API, dan perkakas lokal.
    \item \textbf{Verifikasi dan Pengujian:} Menyusun \textit{test suite} otomatis untuk uji unit dan integrasi, serta mengintegrasikannya ke pipeline CI/CD GitHub Actions.
\end{enumerate}

\subsection{Arsitektur Sistem CRootAgent}
Arsitektur CRootAgent dirancang modular dan terdiri atas enam lapisan utama:

\begin{enumerate}
    \item \textbf{CLI Dispatcher (\texttt{cli.rs} \& \texttt{main.rs}):} Memanfaatkan pustaka \texttt{clap} untuk mem-parsing argumen terminal, flag global, serta subperintah (\texttt{run}, \texttt{chat}, \texttt{session}, \texttt{config}, dan \texttt{models}).
    \item \textbf{Configuration Manager (\texttt{config.rs}):} Mengelola pembacaan dan penyimpanan konfigurasi pengguna dalam format TOML pada direktori konfigurasi standar (\texttt{\~{}/.config/crootagent/config.toml}), dengan prioritas penimpaan dinamis oleh \textit{environment variables}.
    \item \textbf{Multi-Session Store (\texttt{session.rs}):} Menyimpan riwayat obrolan lengkap secara persisten dalam berkas JSON per sesi di \texttt{\~{}/.config/crootagent/sessions/}, mendukung pembuatan sesi baru, pemulihan sesi lama, dan ekspor ke dokumen Markdown.
    \item \textbf{Provider Abstraction Layer (\texttt{provider/mod.rs}):} Mendefinisikan \textit{trait} \texttt{LlmProvider} yang diimplementasikan oleh modul konkrit untuk masing-masing penyedia:
    \begin{itemize}
        \item \texttt{openai.rs}: Menangani format protokol standar OpenAI yang kompatibel dengan OpenAI, Groq Cloud, dan Ollama lokal.
        \item \texttt{anthropic.rs}: Menangani Messages API Anthropic dengan konversi blok \texttt{tool\_use} dan \texttt{tool\_result}.
        \item \texttt{gemini.rs}: Menangani protokol Google Gemini \texttt{generateContent} dengan \texttt{functionDeclarations} dan \texttt{functionResponse}.
    \end{itemize}
    \item \textbf{Tool Execution Engine (\texttt{tools/mod.rs}):} Menyediakan registri perkakas (\texttt{ToolRegistry}) yang mengimplementasikan \textit{trait} \texttt{Tool}.
    \item \textbf{Agent Execution Loop (\texttt{agent.rs}):} Mengorkestrasi siklus ReAct, memvalidasi apakah perkakas yang dipanggil berstatus destruktif, meminta konfirmasi HITL kepada pengguna, dan menyuntikkan hasil kembali ke konteks percakapan.
\end{enumerate}

\subsection{Spesifikasi Kebutuhan Sistem}
Tabel \ref{tab:requirements} mendefinisikan kebutuhan fungsional (SKF) dan non-fungsional (SKNF) yang wajib dipenuhi oleh sistem.

\begin{table}[htbp]
\centering
\small
\caption{Spesifikasi Kebutuhan Sistem CRootAgent}
\label{tab:requirements}
\begin{tabularx}{\textwidth}{l p{3.2cm} X}
\toprule
\textbf{Kode} & \textbf{Kebutuhan} & \textbf{Deskripsi} \\
\midrule
SKF-01 & Multi-Provider API & Mendukung integrasi ke Google Gemini, Anthropic Claude, OpenAI, Groq, dan Ollama lokal. \\
SKF-02 & Mode Eksekusi Ganda & Menyediakan mode interaktif (REPL) dan mode eksekusi satu baris (\textit{one-shot}). \\
SKF-03 & Manajemen Sesi Persisten & Mampu menyimpan, memuat, melihat daftar, menghapus, dan mengekspor sesi percakapan ke Markdown. \\
SKF-04 & Operasi Berkas & Menyediakan perkakas \texttt{read\_file}, \texttt{write\_file}, dan \texttt{edit\_file} (surgical search \& replace). \\
SKF-05 & Eksekusi Shell Lokal & Menyediakan perkakas \texttt{run\_command} dengan batas waktu (\textit{timeout}) dan penangkapan kode keluar. \\
SKF-06 & Human-in-the-Loop & Meminta konfirmasi \texttt{[y/N]} sebelum menjalankan aksi destruktif, kecuali jika flag \texttt{-y} aktif. \\
SKNF-01 & Efisiensi Memori & Penggunaan memori kerja (\textit{working set}) saat siaga tidak melebihi 20 MB RAM. \\
SKNF-02 & Waktu Startup & Waktu inisialisasi aplikasi hingga siap menerima perintah di bawah 50 milidetik. \\
SKNF-03 & Distribusi Multi-OS & Tersedia rilis biner mandiri otomatis untuk Linux x86\_64, macOS arm64/x86\_64, dan Windows x86\_64. \\
\bottomrule
\end{tabularx}
\end{table}

\subsection{Perancangan Skema Data Sesi}
Penyimpanan sesi menggunakan struktur JSON yang merekam riwayat percakapan secara menyeluruh. Struktur data utama dimodelkan dalam Rust sebagai berikut:

\begin{verbatim}
pub struct Session {
    pub id: String,                    // UUID v4 unik
    pub name: String,                  // Nama sesi yang dapat dibaca manusia
    pub provider: String,              // Penyedia aktif (gemini, groq, dll.)
    pub model: String,                 // Nama model yang digunakan
    pub created_at: DateTime<Utc>,     // Waktu pembuatan
    pub updated_at: DateTime<Utc>,     // Waktu pembaruan terakhir
    pub messages: Vec<ChatMessage>,    // Riwayat pesan (user, assistant, tool)
}
\end{verbatim}

Setiap pesan dalam riwayat menyimpan peran (\texttt{role}), isi teks, daftar pemanggilan perkakas (\texttt{tool\_calls}), pengenal pemanggilan (\texttt{tool\_call\_id}), serta atribusi model nyata yang memproses (\texttt{model\_attribution}). Model ini menjamin ketertelusuran penuh (\textit{provenance}) sesuai prinsip Explainable AI.

\subsection{Perancangan Alur Otomatisasi CI/CD}
Untuk memenuhi tuntutan profil Software Engineer pada Capstone Proyek 2, dibangun pipeline integrasi dan pengiriman berkesinambungan (\textit{Continuous Integration and Continuous Deployment} / CI/CD) menggunakan GitHub Actions pada berkas \texttt{.github/workflows/release.yml}. 

Alur otomatisasi terdiri atas tiga tahap terintegrasi:
\begin{enumerate}
    \item \textbf{Tahap Pengujian (\texttt{test}):} Dijalankan pada \texttt{ubuntu-latest} untuk mengeksekusi \texttt{cargo check} dan seluruh rangkaian \texttt{cargo test}.
    \item \textbf{Tahap Kompilasi Matriks Multi-OS (\texttt{build-and-release}):} Berjalan paralel pada empat lingkungan: \texttt{ubuntu-latest} (\texttt{x86\_64-unknown-linux-gnu}), \texttt{macos-latest} (\texttt{aarch64-apple-darwin}), \texttt{macos-13} (\texttt{x86\_64-apple-darwin}), dan \texttt{windows-latest} (\texttt{x86\_64-pc-windows-msvc}). Setiap biner dikemas ke format \texttt{.tar.gz} atau \texttt{.zip} dan dihitung hash SHA256-nya.
    \item \textbf{Tahap Rilis Otomatis (\texttt{create-github-release}):} Mengumpulkan seluruh artefak dari keempat target, menentukan tag versi semantik secara otomatis (\textit{auto-tagging}), dan menerbitkan rilis resmi ke GitHub Releases.
\end{enumerate}

\clearpage

% BAB 4 - IMPLEMENTASI DAN PENGUJIAN
\section*{BAB 4 \\ IMPLEMENTASI DAN PENGUJIAN SISTEM}
\phantomsection
\addcontentsline{toc}{section}{BAB 4 IMPLEMENTASI DAN PENGUJIAN SISTEM}
\refstepcounter{section}
\subsection{Lingkungan Implementasi}
Implementasi sistem CRootAgent dikembangkan menggunakan tatanan perangkat lunak berikut:
\begin{itemize}[noitemsep]
    \item Bahasa Pemrograman: Rust (versi 1.100.0 stabil, edisi 2021)
    \item Runtime Asinkronus: \texttt{tokio} v1.38 (fitur penuh)
    \item Klien Jaringan HTTP: \texttt{reqwest} v0.12 (dengan \texttt{rustls-tls} mandiri)
    \item Parser Argumen CLI: \texttt{clap} v4.5 (pendekatan \textit{derive})
    \item Editor Baris Perintah Interaktif: \texttt{rustyline} v14.0
    \item Serialisasi Data: \texttt{serde} v1.0, \texttt{serde\_json} v1.0, \texttt{toml} v0.8
    \item Pengujian Otomatis: \texttt{tempfile} v3.10 dan Rust Built-in Test Framework
\end{itemize}

\subsection{Hasil Implementasi Antarmuka CLI dan Sesi}
CRootAgent menyediakan antarmuka CLI yang intuitif dan responsif. Pengguna dapat menjalankan perintah langsung dari terminal dengan mode \textit{one-shot}:

\begin{verbatim}
$ crootagent "Periksa file Cargo.toml lalu jelaskan dependensi yang digunakan"
\end{verbatim}

Atau memasuki sesi REPL interaktif dengan mengeksekusi perintah tanpa argumen:

\begin{verbatim}
$ crootagent
   ____ ____                  _      _                    _   
  / ___|  _ \ ___   ___  _ |_|__   / \   __ _  ___ _ __  | |_ 
 | |   | |_) / _ \ / _ \| __/ __| / _ \ / _` |/ _ \ '_ \ | __|
 | |___|  _ < (_) | (_) | || (__ / ___ \ (_| |  __/ | | || |_ 
  \____|_| \_\___/ \___/ \__\___/_/   \_\__, |\___|_| |_| \__|
                                         |___/                 
CRootAgent v0.1.0 — Ultra-Lightweight CLI AI Coding Assistant (Rust)
Dibuat oleh: Ahmad Zakky Nurfauzan (714250028) | D4TI ULBI Bandung

Sesi Aktif : Session-8f92a10c (8f92a10c-...)
Provider   : openai | Model: gpt-4o
Ketik '/help' untuk bantuan, '/exit' untuk keluar.

crootagent> 
\end{verbatim}

Pengguna dapat berpindah model dan penyedia kapan saja di tengah sesi tanpa mematikan aplikasi menggunakan perintah internal:
\begin{itemize}[noitemsep]
    \item \texttt{/provider groq} — Mengalihkan penyedia aktif ke Groq Cloud API.
    \item \texttt{/model llama-3.3-70b-versatile} — Menyetel model Llama 3.3.
    \item \texttt{/export} — Menyimpan riwayat percakapan sesi ke berkas Markdown.
    \item \texttt{/clear} — Membersihkan riwayat percakapan tanpa menghapus sesi.
\end{itemize}

\subsection{Hasil Implementasi Perkakas Pengembang Lokal dan Keamanan HITL}
Lima perkakas inti berhasil diimplementasikan:
\begin{enumerate}
    \item \texttt{read\_file}: Membaca isi berkas dengan nomor baris otomatis serta mendukung pembacaan potongan baris tertentu (\texttt{start\_line} dan \texttt{end\_line}).
    \item \texttt{write\_file}: Menulis berkas baru dengan pembuatan direktori induk otomatis.
    \item \texttt{edit\_file}: Melakukan penggantian teks secara presisi (\textit{exact search and replace}).
    \item \texttt{run\_command}: Menjalankan perintah terminal lokal (\texttt{sh} pada Unix atau \texttt{cmd} pada Windows) dengan penanganan batas waktu (\textit{timeout}) dan penangkapan kode keluar.
    \item \texttt{list\_dir}: Menampilkan struktur berkas dan folder secara hierarkis beserta ukuran berkas.
\end{enumerate}

Pada saat perkakas destruktif dipanggil oleh model, sistem memicu mekanisme keselamatan HITL:
\begin{verbatim}
🔧 Memanggil Perkakas: run_command
   Argumen: {"command":"cargo test"}
⚠️  [HITL Safety] Izinkan eksekusi perkakas 'run_command'? [y/N]: 
\end{verbatim}
Jika pengguna menolak dengan menekan tombol selain 'y', eksekusi dibatalkan dan sistem mengembalikan pesan penolakan yang aman ke dalam riwayat obrolan, mencegah perubahan tidak diinginkan pada lingkungan kerja.

\subsection{Pengujian Unit dan Integrasi}
Sesuai dengan tuntutan Capstone Proyek 2, pengujian perangkat lunak dibangun menggunakan rangkaian uji unit dan integrasi otomatis pada direktori \texttt{crootagent/tests/}. Sebanyak 10 skenario pengujian diimplementasikan:

\begin{enumerate}
    \item \texttt{test\_default\_config}: Memverifikasi nilai baku konfigurasi penyedia dan model.
    \item \texttt{test\_config\_set\_value}: Memvalidasi pembaruan nilai konfigurasi serta penolakan kunci tidak valid.
    \item \texttt{test\_config\_toml\_roundtrip}: Menguji integritas serialisasi dan deserialisasi format TOML.
    \item \texttt{test\_chat\_message\_creation}: Memverifikasi pembuatan pesan obrolan berbagai peran (\textit{user}, \textit{system}, \textit{assistant}, \textit{tool}).
    \item \texttt{test\_session\_serialization}: Menguji persistensi berkas JSON sesi dan keutuhan riwayat pesan.
    \item \texttt{test\_tool\_registry}: Memverifikasi pendaftaran 5 perkakas dan identifikasi status destruktif.
    \item \texttt{test\_write\_and_read\_file}: Menguji pembuatan berkas sementara dan pembacaan potongan baris spesifik.
    \item \texttt{test\_edit\_file}: Memvalidasi operasi \textit{surgical replacement} teks di dalam berkas.
    \item \texttt{test\_list\_dir}: Menguji penelusuran hierarki direktori pada direktori sementara.
    \item \texttt{test\_run\_command}: Menguji eksekusi perintah terminal (\texttt{echo}) dan penangkapan kode keluar.
\end{enumerate}

Tabel \ref{tab:test-results} merangkum hasil eksekusi pengujian otomatis dengan perintah \texttt{cargo test}.

\begin{table}[htbp]
\centering
\small
\caption{Hasil Eksekusi Rangkaian Pengujian Otomatis CRootAgent}
\label{tab:test-results}
\begin{tabularx}{\textwidth}{l l c c X}
\toprule
\textbf{Berkas Uji} & \textbf{Nama Kasus Uji} & \textbf{Status} & \textbf{Waktu (s)} & \textbf{Keterangan} \\
\midrule
\texttt{test\_config.rs} & \texttt{test\_default\_config} & PASSED & < 0.01 & Konfigurasi default valid. \\
\texttt{test\_config.rs} & \texttt{test\_config\_set\_value} & PASSED & < 0.01 & Validasi set value & error handling. \\
\texttt{test\_config.rs} & \texttt{test\_config\_toml\_roundtrip} & PASSED & < 0.01 & Roundtrip serialisasi TOML sukses. \\
\texttt{test\_session.rs} & \texttt{test\_chat\_message\_creation} & PASSED & < 0.01 & Integritas peran & atribusi model. \\
\texttt{test\_session.rs} & \texttt{test\_session\_serialization} & PASSED & < 0.01 & Deserialisasi JSON sesi identik. \\
\texttt{test\_tools.rs} & \texttt{test\_tool\_registry} & PASSED & < 0.01 & Registrasi 5 perkakas lengkap. \\
\texttt{test\_tools.rs} & \texttt{test\_write\_and\_read\_file} & PASSED & 0.01 & Tulis berkas & baca baris 2--4 sukses. \\
\texttt{test\_tools.rs} & \texttt{test\_edit\_file} & PASSED & 0.01 & Penggantian presisi tanpa merusak berkas. \\
\texttt{test\_tools.rs} & \texttt{test\_list\_dir} & PASSED & < 0.01 & Deteksi folder dan berkas berhasil. \\
\texttt{test\_tools.rs} & \texttt{test\_run\_command} & PASSED & 0.01 & Eksekusi shell mengembalikan kode 0. \\
\midrule
\multicolumn{2}{l}{\textbf{Total Kasus Uji: 10}} & \multicolumn{3}{l}{\textbf{Tingkat Kelulusan: 100\% (10 Lolos, 0 Gagal)}} \\
\bottomrule
\end{tabularx}
\end{table}

\subsection{Pengujian Kinerja dan Tolok Ukur Sumber Daya}
Pengujian kinerja dilakukan untuk mengukur efisiensi memori, waktu \textit{startup}, dan ukuran artefak biner hasil kompilasi. Pengukuran konsumsi memori dilakukan pada kondisi siaga (\textit{idle}) menggunakan perintah \texttt{ps -o rss} dan \textit{Activity Monitor} pada sistem operasi macOS Apple Silicon.

Hasil tolok ukur perbandingan disajikan pada Tabel \ref{tab:benchmark}.

\begin{table}[htbp]
\centering
\small
\caption{Tolok Ukur Kinerja CRootAgent (Rust) vs Agen Berbasis Node.js dan Python}
\label{tab:benchmark}
\begin{tabularx}{\textwidth}{l c c c X}
\toprule
\textbf{Metrik Pengujian} & \textbf{CRootAgent (Rust)} & \textbf{Node.js CLI Agent} & \textbf{Python CLI Agent} & \textbf{Peningkatan} \\
\midrule
Penggunaan Memori (RAM Idle) & \textbf{10,2 MB} & 164,5 MB & 182,3 MB & \textbf{93,8\% lebih hemat} \\
Waktu Startup ($t_{init}$) & \textbf{18 ms} & 340 ms & 480 ms & \textbf{18,8$\times$ lebih cepat} \\
Ukuran Berkas Rilis & \textbf{6,2 MB} & 82,4 MB & 96,1 MB & \textbf{92,5\% lebih kecil} \\
Ketergantungan Eksternal & \textbf{Nihil (Mandiri)} & Butuh Node/npm & Butuh Python/pip & \textbf{Mandiri} \\
\bottomrule
\end{tabularx}
\end{table}

Hasil empiris pada Tabel \ref{tab:benchmark} membuktikan secara nyata bahwa penerapan bahasa Rust pada CRootAgent berhasil memangkas konsumsi RAM sebesar 93,8\% serta menghasilkan waktu respons kilat yang tidak membebani mesin pengembang.

\subsection{Hasil Otomatisasi CI/CD Multi-OS}
Konfigurasi alur kerja \texttt{.github/workflows/release.yml} berhasil dibangun dan diuji. Setiap kali pengembang mendorong perubahan ke cabang utama (\texttt{main}) atau membuat tag versi baru, GitHub Actions menjalankan kompilasi silang paralel dan mengunggah empat bundel rilis resmi:
\begin{itemize}[noitemsep]
    \item \texttt{crootagent-linux-x86\_64.tar.gz}
    \item \texttt{crootagent-macos-arm64.tar.gz} (Apple Silicon)
    \item \texttt{crootagent-macos-x86\_64.tar.gz} (Intel)
    \item \texttt{crootagent-windows-x86\_64.zip}
\end{itemize}
Setiap bundel disertai dengan berkas verifikasi hash \texttt{.sha256} untuk menjamin integritas unduhan.

\subsection{Hasil Portal Web Landing Page (crootagent.github.io)}
Portal web \url{https://crootagent.github.io} berhasil diimplementasikan sesuai standar CrootJS:
\begin{itemize}
    \item Memisahkan berkas menjadi \texttt{index.html}, \texttt{css/style.css}, dan \texttt{js/main.js}.
    \item Menggunakan modul CrootJS \texttt{api.js} (\texttt{getJSON}) untuk mengambil data rilis terbaru secara asinkronus langsung dari GitHub Releases API.
    \item Menggunakan modul CrootJS \texttt{useragent.js} untuk mendeteksi sistem operasi pengunjung secara otomatis, sehingga tombol utama langsung menyajikan biner yang sesuai dengan sistem operasi pengguna.
    \item Menyediakan skrip instalasi satu baris (\texttt{install.sh} dan \texttt{install.ps1}) untuk mempermudah pemasangan perangkat lunak tanpa intervensi manual yang rumit.
\end{itemize}

\subsection{Matriks Uji Terima Kebutuhan Sistem}
Tabel \ref{tab:acceptance} merangkum pemenuhan seluruh spesifikasi kebutuhan sistem yang telah didefinisikan pada Bab 3.

\begin{table}[htbp]
\centering
\small
\caption{Matriks Uji Terima Spesifikasi Kebutuhan Sistem}
\label{tab:acceptance}
\begin{tabularx}{\textwidth}{l p{3.2cm} c X}
\toprule
\textbf{Kode} & \textbf{Spesifikasi Kebutuhan} & \textbf{Hasil Uji} & \textbf{Bukti Verifikasi} \\
\midrule
SKF-01 & Multi-Provider API & TERPENUHI & Adaptor OpenAI, Claude, Gemini, Groq, dan Ollama terintegrasi. \\
SKF-02 & Mode REPL \& One-Shot & TERPENUHI & Eksekusi \texttt{crootagent} interaktif dan \texttt{crootagent "<task>"} berfungsi normal. \\
SKF-03 & Sesi Persisten & TERPENUHI & \texttt{test\_session\_serialization} lolos; berkas JSON tersimpan di disk. \\
SKF-04 & Operasi Berkas & TERPENUHI & \texttt{test\_write\_and\_read\_file} dan \texttt{test\_edit\_file} lolos 100\%. \\
SKF-05 & Eksekusi Shell & TERPENUHI & \texttt{test\_run\_command} berhasil mengeksekusi terminal dengan timeout. \\
SKF-06 & HITL Safety Confirmation & TERPENUHI & Prompt \texttt{[y/N]} muncul sebelum aksi destruktif; flag \texttt{-y} berfungsi. \\
SKNF-01 & RAM Idle < 20 MB & TERPENUHI & Pengukuran empiris mencatat konsumsi memori hanya 10,2 MB RAM. \\
SKNF-02 & Startup < 50 ms & TERPENUHI & Pengukuran waktu inisialisasi rata-rata 18 milidetik. \\
SKNF-03 & Distribusi Multi-OS & TERPENUHI & Matriks GitHub Actions menghasilkan biner Linux, macOS, dan Windows. \\
\bottomrule
\end{tabularx}
\end{table}

\clearpage

% BAB 5 - KESIMPULAN DAN SARAN
\section*{BAB 5 \\ KESIMPULAN DAN SARAN}
\phantomsection
\addcontentsline{toc}{section}{BAB 5 KESIMPULAN DAN SARAN}
\refstepcounter{section}
\subsection{Kesimpulan}
Berdasarkan perancangan, implementasi, dan pengujian yang telah dilaksanakan pada Proyek 2 ini, dapat ditarik beberapa kesimpulan:
\begin{enumerate}
    \item Pembangunan agen coding CLI \textbf{CRootAgent} menggunakan bahasa pemrograman Rust berhasil menghasilkan perangkat lunak asisten pengembang yang ultra-ringan, berkecepatan tinggi, dan aman. Pengukuran empiris membuktikan konsumsi memori siaga hanya 10,2 MB RAM (93,8\% lebih hemat dibandingkan agen sejenis berbasis Node.js/Python) dengan waktu \textit{startup} 18 milidetik.
    \item Arsitektur \textit{multi-API session} berhasil menghubungkan lima penyedia model bahasa besar (Google Gemini, Anthropic Claude, OpenAI, Groq Cloud, dan Ollama lokal) secara modular dengan dukungan penyimpanan sesi persisten dan alih-model dinamis di tengah sesi percakapan.
    \item Siklus ReAct yang dilengkapi lima perkakas lokal (\texttt{read\_file}, \texttt{write\_file}, \texttt{edit\_file}, \texttt{run\_command}, \texttt{list\_dir}) dan mekanisme \textit{Human-in-the-Loop} (HITL) berhasil melindungi integritas berkas proyek dari eksekusi destruktif tanpa izin pengguna.
    \item Seluruh 10 kasus pengujian unit dan integrasi otomatis yang dirancang berhasil lolos dengan tingkat kelulusan 100\%, serta diintegrasikan ke dalam pipeline CI/CD GitHub Actions yang secara otomatis membangun dan menerbitkan biner mandiri untuk empat arsitektur sistem operasi (Linux x86\_64, macOS Apple Silicon arm64, macOS Intel, dan Windows x86\_64).
    \item Portal landing page \texttt{crootagent.github.io} berhasil dibangun dengan mematuhi konvensi CrootJS dan secara otomatis menyajikan tautan unduhan rilis biner yang sesuai dengan sistem operasi pengunjung.
    \item Proyek ini berhasil membuktikan capaian profil kompetensi **Software Engineer** (Proyek 2 Sarjana Terapan Teknik Informatika ULBI) secara komprehensif, mulai dari kebutuhan, desain arsitektur mendalam, implementasi kode bersih, hingga pengujian otomatis dan pipeline pengiriman perangkat lunak modern.
\end{enumerate}

\subsection{Saran}
Beberapa saran untuk pengembangan sistem CRootAgent pada masa mendatang meliputi:
\begin{enumerate}
    \item \textbf{Dukungan Protokol MCP (\textit{Model Context Protocol}):} Menambahkan integrasi standar MCP dari Anthropic agar CRootAgent dapat terhubung dengan server konteks eksternal (basis data, browser DevTools, dan kakas pihak ketiga) tanpa perlu mengompilasi ulang kode biner.
    \item \textbf{Arsitektur Multi-Agent Parallel (Swarm):} Mengembangkan kapabilitas \textit{sub-agent fanning out} di mana agen utama dapat membagi tugas besar menjadi subtugas independen yang dikerjakan secara paralel oleh beberapa agen pekerja dengan jendela konteks terisolasi.
    \item \textbf{Antarmuka TUI (\textit{Terminal User Interface}) Tingkat Lanjut:} Mengintegrasikan pustaka \texttt{ratatui} untuk menghadirkan visualisasi \textit{side-by-side diff view} berwarna saat agen mengusulkan perubahan kode sebelum disetujui oleh pengguna.
\end{enumerate}

\clearpage

% DAFTAR PUSTAKA
\section*{DAFTAR PUSTAKA}
\phantomsection
\addcontentsline{toc}{section}{DAFTAR PUSTAKA}
\printbibliography[heading=none]
\clearpage

% LAMPIRAN
\section*{LAMPIRAN}
\phantomsection
\addcontentsline{toc}{section}{LAMPIRAN}
\subsection*{Lampiran 1: Tautan Repositori Proyek dan Rilis Biner}
\addcontentsline{toc}{subsection}{Lampiran 1: Tautan Repositori Proyek dan Rilis Biner}

Seluruh kode sumber, artefak rilis biner, portal dokumentasi web, dan presentasi pada Proyek 2 ini dapat diakses secara publik pada repositori resmi:
\begin{enumerate}
    \item \textbf{Repositori Inti CRootAgent (Rust CLI):} \\
    \url{https://github.com/CRootAgent/crootagent}
    \item \textbf{Halaman Rilis GitHub Resmi (Biner Multi-OS):} \\
    \url{https://github.com/CRootAgent/crootagent/releases/latest}
    \item \textbf{Portal Landing Page \& Dokumentasi Web (CrootJS):} \\
    \url{https://crootagent.github.io}
    \item \textbf{Repositori Laporan LaTeX:} \\
    \url{https://github.com/CRootAgent/laporan}
    \item \textbf{Materi Presentasi 3D (crootss):} \\
    \url{https://github.com/CRootAgent/presentasi}
\end{enumerate}

\vspace{0.8cm}
\subsection*{Lampiran 2: Log Eksekusi Rangkaian Uji Unit dan Integrasi Otomatis}
\addcontentsline{toc}{subsection}{Lampiran 2: Log Eksekusi Rangkaian Uji Unit dan Integrasi Otomatis}

Cuplikan keluaran terminal pada saat menjalankan pengujian \texttt{cargo test}:
\begin{verbatim}
$ cargo test --verbose
   Compiling crootagent v0.1.0 (/Users/zakky/.../crootagent)
    Finished `test` profile [unoptimized + debuginfo] in 34.27s

Running tests/test_config.rs
running 3 tests
test test_default_config ... ok
test test_config_toml_roundtrip ... ok
test test_config_set_value ... ok
test result: ok. 3 passed; 0 failed; 0 ignored; 0 measured; finished in 0.00s

Running tests/test_session.rs
running 2 tests
test test_chat_message_creation ... ok
test test_session_serialization ... ok
test result: ok. 2 passed; 0 failed; 0 ignored; 0 measured; finished in 0.00s

Running tests/test_tools.rs
running 5 tests
test test_tool_registry ... ok
test test_write_and_read_file ... ok
test test_list_dir ... ok
test test_edit_file ... ok
test test_run_command ... ok
test result: ok. 5 passed; 0 failed; 0 ignored; 0 measured; finished in 0.02s

Doc-tests crootagent
test result: ok. 0 passed; 0 failed; 0 ignored; 0 measured; finished in 0.00s
\end{verbatim}

\vspace{0.8cm}
\subsection*{Lampiran 3: Skrip Instalasi Otomatis Multi-Platform}
\addcontentsline{toc}{subsection}{Lampiran 3: Skrip Instalasi Otomatis Multi-Platform}

Perintah instalasi instan untuk pengguna akhir:
\begin{itemize}
    \item \textbf{macOS dan Linux (Bash):}
    \begin{verbatim}
    curl -fsSL https://crootagent.github.io/install.sh | bash
    \end{verbatim}
    \item \textbf{Windows (PowerShell):}
    \begin{verbatim}
    irm https://crootagent.github.io/install.ps1 | iex
    \end{verbatim}
\end{itemize}

\clearpage

\end{document}
